\documentclass[10pt]{article}

% -------------------------
% PAQUETES
% -------------------------
\usepackage[spanish]{babel}
\usepackage[utf8]{inputenc}
\usepackage[T1]{fontenc}
\usepackage{geometry}
\usepackage{setspace}
\usepackage{graphicx}
\usepackage{amsmath}
\usepackage{booktabs}
\usepackage{caption}
\usepackage{hyperref}
\usepackage{apacite}

% -------------------------
% CONFIGURACIÓN DE PÁGINA (según pauta)
% -------------------------
\geometry{
 top=1.5cm,
 left=1.5cm,
 right=1.5cm,
 bottom=2cm
}

\setstretch{1.0} % interlineado simple
\setlength{\parindent}{0pt}
\setlength{\parskip}{0.5em}

% -------------------------
% CAPTIONS (Fig. N. / Tabla N. debajo)
% -------------------------
\captionsetup{font=small, labelfont=bf}
\renewcommand{\figurename}{Fig.}
\renewcommand{\tablename}{Tabla}

% -------------------------
% DATOS (Identificación superior, no ocupa hoja)
% -------------------------
\title{\textbf{Título del Trabajo / Control}}
\author{Natalay Huaiquinao, Carlos Ulloa\\
\texttt{nhuaiquinao2021@alu.uct.cl, carlos.ulloa2020@alu.uct.cl}}
\date{\vspace{-0.6em}Control 1, INFO1184, Semestre I-2026}

% -------------------------
\begin{document}

\maketitle
\vspace{-1.0cm}

% -------------------------
% 2) RESUMEN (100 a 130 palabras)
% -------------------------
\begin{abstract}
Este resumen debe contener entre 100 y 130 palabras. Incluya el objetivo del control, la metodología o estrategia utilizada para abordar el problema (por ejemplo: análisis de requerimientos, diseño de solución paso a paso, validación con casos de prueba), los resultados principales obtenidos y la conclusión más relevante. Debe ser claro y directo, permitiendo comprender el aporte del trabajo sin leer el documento completo. Evite definiciones extensas y detalles secundarios; priorice qué se hizo, cómo se hizo y qué se obtuvo. Si corresponde, mencione brevemente el alcance y limitaciones del desarrollo realizado.
\end{abstract}

% -------------------------
% 3) INTRODUCCIÓN
% -------------------------
\section{Introducción}
Aquí se presenta el contexto del problema y la problemática abordada. Incluya (según corresponda) preguntas de investigación, caso de estudio o enunciado del ejercicio. Indique objetivos (general y específicos) y describa brevemente la estructura del informe (Introducción, Desarrollo, Reflexión, Referencias y Anexo).

% -------------------------
% 4) DESARROLLO
% -------------------------
\section{Desarrollo}
Desarrolle la solución paso a paso, integrando contenidos vistos en clases y, si aplica, lecturas y dos fuentes de referencia. Explique supuestos, decisiones y validación con ejemplos.

\subsection{Ejemplo de figura (numeración consecutiva y rótulo inferior)}
\begin{figure}[h]
    \centering
    % Reemplace ejemplo.png por su imagen real
    \includegraphics[width=0.8\textwidth]{ejemplo.png}
    \caption{Título de la figura. (\citeauthor{autorEjemplo2020}, \citeyear{autorEjemplo2020})}
    \label{fig:ejemplo}
\end{figure}

\subsection{Ejemplo de tabla (numeración consecutiva y rótulo inferior)}
\begin{table}[h]
\centering
\begin{tabular}{ccc}
\toprule
Columna 1 & Columna 2 & Columna 3 \\
\midrule
Dato 1 & Dato 2 & Dato 3 \\
\bottomrule
\end{tabular}
\caption{Título de la tabla. (\citeauthor{autorEjemplo2020}, \citeyear{autorEjemplo2020})}
\label{tab:ejemplo}
\end{table}

\subsection{Ejemplo de ecuación (numeración a la derecha)}
La ecuación debe ir numerada automáticamente. Si proviene de una fuente, menciónela en el texto.
\begin{equation}
E = mc^2
\label{eq:energia}
\end{equation}

% -------------------------
% 5) REFLEXIÓN
% -------------------------
\section{Reflexión}
Justifique las decisiones tomadas (por ejemplo: elección de metodología, estructura de la solución, herramientas utilizadas, criterios de validación). Analice resultados, dificultades, mejoras posibles y una reflexión crítica del proceso (qué funcionó, qué no, y cómo lo abordaría mejor en una siguiente iteración).

% -------------------------
% 6) REFERENCIAS (APA, solo dos)
% -------------------------
\section{Referencias}
\bibliographystyle{apacite}
\bibliography{referencias}

% -------------------------
% 7) ANEXO (no cuenta para 2-3 hojas del informe principal)
% -------------------------
\appendix
\section{Anexo}
Use esta sección solo si necesita profundizar procedimientos o agregar evidencias parciales. Si existe código, la pauta sugiere NO incluirlo en el cuerpo principal; aquí puede agregar solo partes importantes o procedimientos relevantes.

\subsection{Bitácora de trabajo}
\begin{itemize}
    \item Día 1 (AAAA-MM-DD): Descripción breve de actividades realizadas, acuerdos, decisiones y avances.
    \item Día 2 (AAAA-MM-DD): Avances, problemas encontrados, cómo se resolvieron o próximos pasos.
\end{itemize}

\end{document}